\documentclass[12pt,oneside,a4paper]{article}

\usepackage[margin=1in]{geometry}
\usepackage{graphicx}
\usepackage{hyperref}
\usepackage{xcolor}
\usepackage{fancyhdr}
\usepackage{listings}
\usepackage{tabularx}
\usepackage{float}
\usepackage{titlesec}
\usepackage{textcomp}

% Set up page style
\pagestyle{fancy}
\fancyhf{}
\rhead{UIT-SE357 LMS}
\lhead{Software Architecture Document}
\cfoot{\thepage}

% Code formatting
\lstset{
    basicstyle=\ttfamily\small,
    breaklines=true,
    columns=fullflexible,
    frame=single,
    language=json,
    keywordstyle=\color{blue},
    commentstyle=\color{gray},
    stringstyle=\color{red},
}

% Title formatting
\titleformat{\section}{\Large\bfseries}{\thesection}{1em}{}
\titleformat{\subsection}{\large\bfseries}{\thesubsection}{1em}{}

\title{\textbf{Software Architecture Document}\\
UIT-SE357: Learning Management System}
\author{Le Minh Phat \& Team}
\date{May 15, 2026}

\begin{document}

% ===== TITLE PAGE =====
\begin{titlepage}
    \begin{center}
        \vspace*{2cm}
        {\LARGE \textbf{SOFTWARE ARCHITECTURE DOCUMENT}}\\
        \vspace{0.5cm}
        {\large UIT-SE357: Learning Management System}\\
        \vspace{3cm}
        {\large Version 1.0}\\
        \vspace{0.5cm}
        {\large May 15, 2026}\\
        \vspace{3cm}
        Document Status: \textbf{Final}\\
        \vspace{2cm}
        \small
        \begin{tabularx}{\textwidth}{lX}
            \textbf{Prepared by:} & Le Minh Phat \\
            \textbf{Reviewed by:} & UIT-SE357 Team \\
            \textbf{Approved by:} & Project Manager \\
        \end{tabularx}
        \vfill
    \end{center}
\end{titlepage}

% ===== DOCUMENT CONTROL =====
\section*{Document Control}

\subsection*{Distribution}
This document is intended for stakeholders, architects, developers, and technical reviewers of the UIT-SE357 Learning Management System project.

\subsection*{Confidentiality}
This document contains technical information about the system architecture and is available to project team members.

\subsection*{Version History}

\begin{tabularx}{\textwidth}{|p{1.5cm}|p{2cm}|p{3cm}|X|}
    \hline
    \textbf{Version} & \textbf{Date} & \textbf{Author} & \textbf{Changes} \\
    \hline
    1.0 & 2026-05-15 & LeMinhPhat & Initial SAD document covering implemented ASRs \\
    \hline
\end{tabularx}

\newpage

% ===== TABLE OF CONTENTS =====
\tableofcontents
\newpage

% ===== SECTION 1: INTRODUCTION =====
\section{Introduction}

\subsection{Purpose}
This Software Architecture Document (SAD) provides a comprehensive description of the implemented architecture for the UIT-SE357 Learning Management System. It documents the architectural decisions, design patterns, and component interactions that support the system's key attributes and requirements.

\subsection{Scope}
This document covers the implemented Architectural Significant Requirements (ASRs) within the LMS, including:
\begin{itemize}
    \item Authentication and Authorization (ASR-SEC-02, ASR-SEC-03, ASR-ENROLL-01)
    \item Audit Logging and Compliance (ASR-SEC-11)
    \item Performance Optimization (ASR-PERF-01, ASR-PERF-03)
    \item Availability and Reliability (ASR-AVAIL-01)
\end{itemize}

This document does not cover unimplemented ASRs or planned future capabilities.

\subsection{Audience}
\begin{itemize}
    \item Software Architects and Senior Developers
    \item DevOps and Infrastructure Teams
    \item Security and Compliance Officers
    \item Project Managers and Stakeholders
\end{itemize}

\newpage

% ===== SECTION 2: SYSTEM OVERVIEW =====
\section{System Overview}

\subsection{System Context}
The UIT-SE357 Learning Management System is a web-based platform designed to support educational institutions in managing courses, assignments, materials, and student evaluations. The system serves three primary user roles: students, instructors, and administrators.

\subsection{Technology Stack}

\begin{tabularx}{\textwidth}{|l|X|}
    \hline
    \textbf{Component} & \textbf{Technology} \\
    \hline
    Frontend & React.js + TypeScript + Vite \\
    Backend & Express.js + Node.js + TypeScript \\
    Database & PostgreSQL 16 \\
    Cache Store & Redis 7 \\
    Reverse Proxy & Caddy v2 \\
    Container Runtime & Docker + Docker Compose \\
    \hline
\end{tabularx}

\subsection{Deployment Architecture}
The system is deployed using Docker Compose with the following components:
\begin{itemize}
    \item \textbf{Frontend (web)}: React SPA served via Nginx
    \item \textbf{Backend Primary}: Express.js server instance
    \item \textbf{Backend Secondary}: Express.js server instance (failover)
    \item \textbf{Database}: PostgreSQL with persistent volume
    \item \textbf{Cache}: Redis for session \& rate-limit storage
    \item \textbf{Reverse Proxy}: Caddy for load balancing \& TLS termination
\end{itemize}

\newpage

% ===== SECTION 3: ARCHITECTURAL OVERVIEW =====
\section{Architectural Overview}

\subsection{High-Level Architecture}
The LMS follows a layered architecture pattern:

\begin{lstlisting}
┌─────────────────────────────────────┐
│   Frontend (React SPA)               │
└──────────────┬──────────────────────┘
               │ HTTP/HTTPS (TLS)
┌──────────────▼──────────────────────┐
│   Caddy Reverse Proxy (Load Balancer)│
│   - Health checks                    │
│   - TLS termination                  │
└──┬────────────────────────────────┬──┘
   │                                │
┌──▼──────────┐            ┌──────▼──┐
│Backend      │            │Backend   │
│Primary      │            │Secondary │
└──┬──────────┘            └──────┬───┘
   │                              │
   └──────────────┬───────────────┘
                  │ Shared State
        ┌─────────▼────────────┐
        │ PostgreSQL Database  │
        │ Redis Cache          │
        └──────────────────────┘
\end{lstlisting}

\subsection{Key Design Principles}
\begin{itemize}
    \item \textbf{Separation of Concerns}: Clear separation between authentication, business logic, and data access
    \item \textbf{Defense in Depth}: Multiple security layers (authentication, authorization, rate limiting, audit logging)
    \item \textbf{High Availability}: Redundant backend instances with automatic health-check based failover
    \item \textbf{Performance}: Redis caching, response compression, and rate limiting
    \item \textbf{Auditability}: Comprehensive audit logging for compliance and troubleshooting
\end{itemize}

\newpage

% ===== SECTION 4: IMPLEMENTED ASRS =====
\section{Implemented Architectural Significant Requirements}

\subsection{ASR-SEC-02: JWT Authentication with Refresh Tokens}

\subsubsection{Overview}
This ASR implements stateless authentication using JSON Web Tokens (JWT) with short-lived access tokens and refresh tokens.

\subsubsection{Components}
\begin{itemize}
    \item \textbf{Auth Controller} (\texttt{auth/auth.controller.ts})
        \begin{itemize}
            \item Handles POST /auth/login, /auth/register, /auth/refresh endpoints
            \item Validates input and returns JWT tokens
        \end{itemize}
    \item \textbf{Auth Service} (\texttt{auth/auth.service.ts})
        \begin{itemize}
            \item User registration with password hashing
            \item Login authentication with password verification
            \item Token generation using Node's jsonwebtoken library
        \end{itemize}
    \item \textbf{Auth Middleware} (\texttt{auth/auth.middleware.ts})
        \begin{itemize}
            \item JWT token extraction from Authorization header
            \item Token signature verification
            \item User context injection into request object
        \end{itemize}
\end{itemize}

\subsubsection{Implementation Details}
\begin{lstlisting}
# Environment Configuration
AUTH_SECRET=<32-byte-key>
AUTH_SECRET_EXPIRES_IN=15m
AUTH_REFRESH_SECRET=<32-byte-key>
AUTH_REFRESH_SECRET_EXPIRES_IN=7d
\end{lstlisting}

\subsubsection{Flow}
\begin{enumerate}
    \item User submits email and password via POST /auth/login
    \item Service verifies password against stored hash (bcrypt)
    \item JWT tokens (access + refresh) are generated with configured secrets
    \item Client stores tokens in memory (access) and httpOnly cookie (refresh)
    \item Subsequent requests include access token in Authorization header
    \item Middleware validates token on each protected endpoint
    \item When access token expires, client uses refresh token to obtain new access token
\end{enumerate}

\subsubsection{Security Considerations}
\begin{itemize}
    \item Access tokens are short-lived (15 minutes default)
    \item Refresh tokens use separate, stronger secret and longer expiry (7 days default)
    \item Tokens are signed with HS256 algorithm
    \item Refresh tokens stored in httpOnly cookies (XSS resistant)
    \item Password hashed with bcrypt (salted, adaptive)
\end{itemize}

\subsection{ASR-SEC-03: Login Attempt Lockout}

\subsubsection{Overview}
This ASR protects against brute-force attacks by tracking failed login attempts per IP and implementing a temporary account lockout.

\subsubsection{Components}
\begin{itemize}
    \item \textbf{Login Attempt Tracker} (\texttt{auth/auth.service.ts})
        \begin{itemize}
            \item Queries login attempt history from database
            \item Records each failed attempt with timestamp
            \item Enforces lockout period when threshold exceeded
        \end{itemize}
    \item \textbf{LoginAttempt Table} (PostgreSQL)
        \begin{itemize}
            \item Stores: \texttt{ip}, \texttt{attempts}, \texttt{lastTry}, \texttt{blockedUntil}
            \item Primary key: \texttt{ip}
        \end{itemize}
\end{itemize}

\subsubsection{Configuration}
\begin{lstlisting}
# Environment Configuration
MAX_LOGIN_ATTEMPTS=5
LOCKOUT_TIME_MS=900000  # 15 minutes
\end{lstlisting}

\subsubsection{Algorithm}
\begin{enumerate}
    \item On login attempt, query LoginAttempt table for the client's IP
    \item Check if \texttt{blockedUntil} timestamp is in the future
    \item If blocked, return 429 "Too many attempts"
    \item If password incorrect, increment \texttt{attempts} counter
    \item If \texttt{attempts} >= MAX\_LOGIN\_ATTEMPTS, set \texttt{blockedUntil} to now + LOCKOUT\_TIME\_MS
    \item On successful login, reset \texttt{attempts} to 0
\end{enumerate}

\subsubsection{Security Considerations}
\begin{itemize}
    \item IP-based tracking: Effective for typical scenarios but vulnerable to distributed attacks
    \item Lockout duration: 15 minutes default balances security vs. user convenience
    \item Reset on success: Encourages legitimate users to avoid lockout
    \item Database-backed: Persists across server restarts
\end{itemize}

\subsection{ASR-SEC-11: Audit Log Retention}

\subsubsection{Overview}
This ASR implements comprehensive audit logging with configurable retention policies to support compliance requirements.

\subsubsection{Components}
\begin{itemize}
    \item \textbf{Audit Middleware} (\texttt{middlewares/audit.middleware.ts})
        \begin{itemize}
            \item Intercepts all HTTP requests
            \item Captures: user, action, resource, method, IP, timestamp, response status
        \end{itemize}
    \item \textbf{Audit Logger} (\texttt{util/auditLogger.ts})
        \begin{itemize}
            \item Formats audit events as JSON
            \item Writes to JSONL (JSON Lines) files
            \item Files organized by date: \texttt{audit\_YYYY-MM-DD.jsonl}
        \end{itemize}
    \item \textbf{Audit Query Service} (\texttt{audit/audit.controller.ts})
        \begin{itemize}
            \item GET /api/audit endpoint with filtering
            \item Supports: date range, user ID, action, category
        \end{itemize}
    \item \textbf{Retention Scheduler}
        \begin{itemize}
            \item Cron job to delete logs older than AUDIT\_RETENTION\_DAYS
            \item Runs daily to maintain compliance window
        \end{itemize}
\end{itemize}

\subsubsection{Storage Format}
\begin{lstlisting}
# JSONL Format (one JSON object per line)
{
  "timestamp": "2026-05-15T10:30:45.123Z",
  "user_id": "user-uuid",
  "action": "GET",
  "resource": "/api/courses/course-id",
  "method": "GET",
  "status_code": 200,
  "ip_address": "192.168.1.1",
  "duration_ms": 45
}
\end{lstlisting}

\subsubsection{Configuration}
\begin{lstlisting}
# Environment Configuration
AUDIT_LOG_DIR=/app/logs/audit
AUDIT_RETENTION_DAYS=90
\end{lstlisting}

\subsubsection{Query Interface}
\begin{lstlisting}
GET /api/audit?from=2026-05-01&to=2026-05-15&action=DELETE&user_id=user-uuid&limit=100
\end{lstlisting}

\subsubsection{Security Considerations}
\begin{itemize}
    \item File-based storage: Simple, auditable, no database dependency
    \item JSONL format: Line-based parsing prevents partial read corruption
    \item Configurable retention: Supports regulatory requirements (e.g., 90-day minimum)
    \item Read-only query interface: Logs cannot be modified post-creation
    \item Timestamp immutability: Event timestamp is set at capture time
\end{itemize}

\subsection{ASR-AVAIL-01: Load Balancing with Health Checks}

\subsubsection{Overview}
This ASR provides high availability through redundant backend instances and health-check based failover.

\subsubsection{Components}
\begin{itemize}
    \item \textbf{Caddy Reverse Proxy}
        \begin{itemize}
            \item Listens on :80 (HTTP) and :443 (HTTPS)
            \item Terminates TLS connections
            \item Routes /api/* to backend, / to frontend
        \end{itemize}
    \item \textbf{Load Balancer}
        \begin{itemize}
            \item Policy: 'first' (tries primary, falls back to secondary)
            \item Distributes across backend-primary and backend-secondary
        \end{itemize}
    \item \textbf{Health Check}
        \begin{itemize}
            \item Probes /api/health endpoint every 10 seconds
            \item Marks instance down after 1 failed check
            \item Marks instance up after 1 successful check
        \end{itemize}
    \item \textbf{Health Endpoint} (\texttt{health/health.route.ts})
        \begin{itemize}
            \item Returns 200 OK if system is healthy
            \item Checks: database connectivity, Redis connectivity
        \end{itemize}
\end{itemize}

\subsubsection{Configuration (Caddyfile)}
\begin{lstlisting}
reverse_proxy backend-primary:8000 backend-secondary:8000 {
  lb_policy first
  health_uri /api/health
  health_interval 10s
  health_timeout 3s
  health_fails 1
  health_passes 1
}
\end{lstlisting}

\subsubsection{Failover Behavior}
\begin{itemize}
    \item Primary instance is preferred (first policy)
    \item If primary fails health check, traffic switches to secondary
    \item When primary recovers, traffic returns to primary
    \item Both instances share PostgreSQL database and Redis cache
    \item No session state on instances (stateless design)
\end{itemize}

\subsubsection{Availability Considerations}
\begin{itemize}
    \item RPO (Recovery Point Objective): 0 (no data loss due to stateless design)
    \item RTO (Recovery Time Objective): \textasciitilde 15 seconds (health check interval + detection)
    \item Single point of failure: PostgreSQL database (requires separate replication strategy)
    \item Single point of failure: Caddy proxy (requires separate load balancer or HA setup)
\end{itemize}

\subsection{ASR-PERF-01: Rate Limiting}

\subsubsection{Overview}
This ASR protects the system from resource exhaustion through request rate limiting, backed by Redis for distributed rate limit state.

\subsubsection{Components}
\begin{itemize}
    \item \textbf{Rate Limit Middleware} (\texttt{middlewares/rateLimit.ts})
        \begin{itemize}
            \item Intercepts all requests before route handlers
            \item Uses client IP as identifier (or user ID for authenticated requests)
        \end{itemize}
    \item \textbf{Redis Store}
        \begin{itemize}
            \item Stores: \texttt{client-id:request-count:window}
            \item Enables distributed rate limiting across instances
        \end{itemize}
\end{itemize}

\subsubsection{Configuration}
\begin{lstlisting}
# Environment Configuration
RATE_LIMIT_WINDOW_MS=60000  # 1 minute
RATE_LIMIT_MAX=100          # 100 requests per minute
\end{lstlisting}

\subsubsection{Algorithm}
\begin{enumerate}
    \item Extract client identifier (IP or user ID)
    \item Query Redis for request count in current time window
    \item If count >= RATE_LIMIT_MAX, return 429 "Too Many Requests"
    \item Otherwise, increment counter and set TTL equal to window duration
    \item Proceed with request
\end{enumerate}

\subsubsection{Response Headers}
\begin{lstlisting}
X-RateLimit-Limit: 100
X-RateLimit-Remaining: 87
X-RateLimit-Reset: 1684159845
\end{lstlisting}

\subsubsection{Performance Considerations}
\begin{itemize}
    \item Redis-backed: Adds \textasciitilde 5-10ms latency per request
    \item Sliding window: More accurate than fixed window
    \item Per-IP vs. per-user: IP-based for unauthenticated, user-based for authenticated
    \item Threshold tuning: 100 req/min balances between protection and legitimate usage
\end{itemize}

\subsection{ASR-PERF-03: Response Caching and Compression}

\subsubsection{Overview}
This ASR improves response times through Redis-backed caching and automatic response compression.

\subsubsection{Components}
\begin{itemize}
    \item \textbf{Compression Middleware}
        \begin{itemize}
            \item Automatically gzips/zstd compresses responses
            \item Respects Content-Type for selective compression
        \end{itemize}
    \item \textbf{Cache Middleware}
        \begin{itemize}
            \item Intercepts responses for caching
            \item Cache key: \texttt{method:path:query-string}
        \end{itemize}
    \item \textbf{Redis Cache}
        \begin{itemize}
            \item Stores serialized JSON responses
            \item TTL per resource type (courses: 5min, materials: 5min, etc.)
        \end{itemize}
\end{itemize}

\subsubsection{Configuration}
Cache TTLs are set per endpoint:
\begin{lstlisting}
GET /api/courses       -> TTL: 5 minutes
GET /api/materials     -> TTL: 5 minutes
POST /api/courses      -> Cache invalidation
DELETE /api/courses/:id -> Cache invalidation
\end{lstlisting}

\subsubsection{Compression Ratios}
\begin{itemize}
    \item JSON responses: \textasciitilde 70-80\% compression ratio
    \item HTML: \textasciitilde 80-85\% compression ratio
    \item Uncompressed response: \textasciitilde 50-100KB for typical course list
    \item Compressed response: \textasciitilde 10-20KB
\end{itemize}

\subsubsection{Cache Invalidation Strategy}
\begin{itemize}
    \item Explicit invalidation: POST/PUT/DELETE operations clear related cache entries
    \item TTL-based expiration: Automatic expiration after configured duration
    \item No cache stamping: Simple invalidation (could add stampeding prevention if needed)
\end{itemize}

\subsection{ASR-ENROLL-01: Role-Based Access Control}

\subsubsection{Overview}
This ASR implements role-based access control to protect sensitive resources and enforce enrollment requirements.

\subsubsection{Roles}
\begin{itemize}
    \item \textbf{STUDENT}: Can view own materials, submit assignments, view own feedback
    \item \textbf{INSTRUCTOR}: Can manage materials, grade submissions, view all enrolled students
    \item \textbf{ADMIN}: Full system access (user management, configuration, audit logs)
\end{itemize}

\subsubsection{Components}
\begin{itemize}
    \item \textbf{Role Extractor}
        \begin{itemize}
            \item Extracts role from JWT token (or database lookup)
            \item Stores in request context
        \end{itemize}
    \item \textbf{Enrollment Checker}
        \begin{itemize}
            \item Verifies user is enrolled in requested class
            \item Checks Enrollment table: \texttt{user\_id, class\_id, role}
        \end{itemize}
    \item \textbf{RBAC Policy Engine}
        \begin{itemize}
            \item Evaluates permissions: (user\_role, resource, action)
            \item Returns allow/deny decision
        \end{itemize}
\end{itemize}

\subsubsection{Access Control Matrix}

\begin{tabularx}{\textwidth}{|l|c|c|c|}
    \hline
    \textbf{Resource} & \textbf{Student} & \textbf{Instructor} & \textbf{Admin} \\
    \hline
    View course materials & Y (enrolled) & Y & Y \\
    Create materials & N & Y & Y \\
    Grade submissions & N & Y & Y \\
    View audit logs & N & N & Y \\
    Manage users & N & N & Y \\
    \hline
\end{tabularx}

\subsubsection{Implementation Locations}
\begin{itemize}
    \item \texttt{materials/materials.service.ts}: Enrollment check before returning materials
    \item \texttt{submissions/submissions.service.ts}: RBAC check for viewing/grading
    \item \texttt{feedback/feedback.service.ts}: Enrollment verification
    \item \texttt{audit/audit.controller.ts}: Admin-only check
\end{itemize}

\newpage

% ===== SECTION 5: DEPLOYMENT =====
\section{Deployment Architecture}

\subsection{Development Environment}
\begin{lstlisting}
docker-compose up -d
# Services available at:
# Frontend: http://localhost:5173
# Backend: http://localhost:8000
# Database: localhost:5432
# Redis: localhost:6379
\end{lstlisting}

\subsection{Production Deployment}
\begin{itemize}
    \item Docker Compose for container orchestration
    \item Two backend instances for redundancy
    \item Shared PostgreSQL database
    \item Shared Redis cache
    \item Caddy reverse proxy with TLS
    \item Log volume mount at /app/logs for audit logs
\end{itemize}

\subsection{Database Schema}
Key tables:
\begin{itemize}
    \item \texttt{User}: user\_id, email, password\_hash, role, created\_at
    \item \texttt{LoginAttempt}: ip, attempts, lastTry, blockedUntil
    \item \texttt{Enrollment}: user\_id, class\_id, role, enrolled\_at
    \item \texttt{Material}: id, class\_id, title, content, created\_at
    \item \texttt{Submission}: id, assignment\_id, user\_id, content, submitted\_at
    \item \texttt{Feedback}: id, submission\_id, user\_id, comment, created\_at
\end{itemize}

\newpage

% ===== SECTION 6: QUALITY ATTRIBUTES =====
\section{Quality Attributes}

\subsection{Security}
\begin{itemize}
    \item \textbf{Authentication}: JWT-based, stateless, XSS/CSRF resistant
    \item \textbf{Authorization}: Role-based with enrollment validation
    \item \textbf{Audit}: Comprehensive logging of all user actions
    \item \textbf{Rate Limiting}: Protects against denial-of-service attacks
    \item \textbf{Password Security}: Bcrypt hashing with random salt
\end{itemize}

\subsection{Availability}
\begin{itemize}
    \item \textbf{RTO}: \textasciitilde 15 seconds (automatic failover via health checks)
    \item \textbf{RPO}: 0 (stateless design with shared database)
    \item \textbf{Uptime}: Target 99.5\% (with single points of failure)
    \item \textbf{Health Checks}: Every 10 seconds on all instances
\end{itemize}

\subsection{Performance}
\begin{itemize}
    \item \textbf{Response Compression}: 70-80\% reduction in bandwidth
    \item \textbf{Caching}: 5-minute TTL on read-heavy endpoints
    \item \textbf{Rate Limiting}: 100 req/min per client
    \item \textbf{Load Balancing}: Distributes traffic across instances
    \item \textbf{Database Indexing}: Indexes on common query fields
\end{itemize}

\subsection{Maintainability}
\begin{itemize}
    \item \textbf{Code Organization}: Modular service structure (auth, courses, materials, etc.)
    \item \textbf{Logging}: Structured logging for debugging and monitoring
    \item \textbf{Documentation}: SAD documents architectural decisions
    \item \textbf{Configuration}: Environment-based configuration (no hardcoded secrets)
\end{itemize}

\newpage

% ===== SECTION 7: UNIMPLEMENTED ASRS =====
\section{Unimplemented ASRs}

The following ASRs are defined in the requirements but not yet implemented:

\begin{itemize}
    \item \textbf{ASR-AGGR-01}: Data aggregation and reporting
    \item \textbf{ASR-AVAIL-02}: Database replication and backup strategy
    \item \textbf{ASR-PERF-02}: Database query optimization and indexing strategy
    \item \textbf{ASR-NOTIF-01}: Real-time notifications and event streaming
    \item \textbf{ASR-FILE-01}: File storage and CDN integration (materials handled in DB for now)
    \item \textbf{ASR-SEC-01}: TLS/HTTPS encryption (partially implemented via Caddy)
    \item \textbf{ASR-SEC-04 through ASR-SEC-10}: Additional security requirements
    \item \textbf{ASR-SEC-12 through ASR-SEC-15}: Additional compliance requirements
\end{itemize}

These will be addressed in future phases of development.

\newpage

% ===== SECTION 8: RISKS AND MITIGATIONS =====
\section{Risks and Mitigations}

\subsection{Single Point of Failure: Database}
\textbf{Risk}: PostgreSQL instance is not replicated. Database failure causes complete system downtime.
\newline\textbf{Mitigation}: Implement PostgreSQL streaming replication or managed database service with automatic failover.

\subsection{Single Point of Failure: Caddy Load Balancer}
\textbf{Risk}: Caddy instance failure removes request routing capability.
\newline\textbf{Mitigation}: Deploy Caddy in HA configuration with Keepalived or similar VIP manager.

\subsection{Audit Log Storage on Shared Volume}
\textbf{Risk}: Audit logs stored on shared volume; if volume fails, audit trail is lost.
\newline\textbf{Mitigation}: Implement audit log replication to dedicated storage or centralized logging service.

\subsection{Rate Limiting via IP}
\textbf{Risk}: IP-based rate limiting vulnerable to distributed attacks and proxy scenarios.
\newline\textbf{Mitigation}: Supplement with user-based rate limiting for authenticated requests.

\subsection{JWT Secret Rotation}
\textbf{Risk}: JWT secrets stored in environment variables; rotation requires deployment.
\newline\textbf{Mitigation}: Implement key rotation mechanism with algorithm versioning in JWT header.

\newpage

% ===== SECTION 9: RECOMMENDATIONS =====
\section{Recommendations}

\subsection{Short-term (1-3 months)}
\begin{enumerate}
    \item Implement database backup and restore procedures
    \item Add centralized logging aggregation (ELK stack or CloudWatch)
    \item Implement security scanning in CI/CD pipeline
    \item Add integration tests for audit logging
    \item Document operational runbooks for common scenarios
\end{enumerate}

\subsection{Medium-term (3-6 months)}
\begin{enumerate}
    \item Implement database replication for high availability
    \item Migrate to managed database service (AWS RDS, Azure Database)
    \item Implement distributed tracing for performance analysis
    \item Add automated performance testing and benchmarking
    \item Implement API versioning strategy
\end{enumerate}

\subsection{Long-term (6-12 months)}
\begin{enumerate}
    \item Migrate to Kubernetes for container orchestration
    \item Implement service mesh (Istio) for advanced traffic management
    \item Add machine learning-based anomaly detection for security
    \item Implement advanced caching strategies (cache warming, predictive)
    \item Multi-region deployment with geo-failover
\end{enumerate}

\newpage

% ===== SECTION 10: REFERENCES =====
\section{References}

\subsection{External Standards}
\begin{itemize}
    \item RFC 7519: JSON Web Token (JWT) Specification
    \item OWASP Top 10: Web Application Security Risks
    \item NIST Cybersecurity Framework
    \item ISO/IEC 27001: Information Security Management
\end{itemize}

\subsection{Project Documentation}
\begin{itemize}
    \item \texttt{README.md}: Quick start guide
    \item \texttt{docker-compose.yml}: Deployment configuration
    \item \texttt{.env.example}: Environment variables
    \item \texttt{Caddyfile}: Load balancer configuration
\end{itemize}

\subsection{Architecture Diagrams}
\begin{itemize}
    \item ASR-SEC-02-L4.puml: JWT Authentication (C4 Level 4)
    \item ASR-SEC-03-L4.puml: Login Lockout (C4 Level 4)
    \item ASR-SEC-11-L4.puml: Audit Logging (C4 Level 4)
    \item ASR-AVAIL-01-L4.puml: Load Balancing (C4 Level 4)
    \item ASR-PERF-01-L4.puml: Rate Limiting (C4 Level 4)
    \item ASR-PERF-03-L4.puml: Caching \& Compression (C4 Level 4)
    \item ASR-ENROLL-01-L4.puml: Role-Based Access Control (C4 Level 4)
\end{itemize}

\newpage

% ===== APPENDIX A =====
\section*{Appendix A: Configuration Reference}

\subsection*{Environment Variables}

\begin{lstlisting}
# Application
NODE_ENV=production
APP_HOST=0.0.0.0
APP_PORT=8000

# Database
DATABASE_URL=postgresql://user:pass@db:5432/uit_se357

# Authentication
AUTH_SECRET=<32-byte-random-string>
AUTH_SECRET_EXPIRES_IN=15m
AUTH_REFRESH_SECRET=<32-byte-random-string>
AUTH_REFRESH_SECRET_EXPIRES_IN=7d

# Security
MAX_LOGIN_ATTEMPTS=5
LOCKOUT_TIME_MS=900000

# Performance
RATE_LIMIT_WINDOW_MS=60000
RATE_LIMIT_MAX=100

# Audit
AUDIT_LOG_DIR=/app/logs/audit
AUDIT_RETENTION_DAYS=90

# Cache
REDIS_HOST=redis
REDIS_PORT=6379

# Email (optional)
EMAIL_HOST=smtp.gmail.com
EMAIL_PORT=587
EMAIL_USER=your-email@gmail.com
EMAIL_PASSWORD=your-app-password
\end{lstlisting}

\newpage

% ===== APPENDIX B =====
\section*{Appendix B: Component File Mapping}

\begin{tabularx}{\textwidth}{|l|X|}
    \hline
    \textbf{Component} & \textbf{File Path} \\
    \hline
    Auth Controller & \texttt{server/src/auth/auth.controller.ts} \\
    Auth Service & \texttt{server/src/auth/auth.service.ts} \\
    Auth Middleware & \texttt{server/src/auth/auth.middleware.ts} \\
    Auth Config & \texttt{server/src/config/auth.config.ts} \\
    Audit Middleware & \texttt{server/src/middlewares/audit.middleware.ts} \\
    Audit Logger & \texttt{server/src/util/auditLogger.ts} \\
    Audit Controller & \texttt{server/src/audit/audit.controller.ts} \\
    Rate Limit Middleware & \texttt{server/src/middlewares/rateLimit.ts} \\
    Health Route & \texttt{server/src/health/health.route.ts} \\
    Compression Middleware & \texttt{server/src/app.ts} (compression package) \\
    Materials Service & \texttt{server/src/materials/materials.service.ts} \\
    \hline
\end{tabularx}

\end{document}
